\chapter{Project Overview and Boundaries}
\label{ch:orientation}
\chapteroverview{Establish the project's purpose, supported surfaces, explicit
boundaries, and evidence hierarchy before any implementation detail is
introduced.}{Onboarding to the repository, deciding whether a requested use is
supported, or locating the authoritative source for a high-level claim.}

\section{Chapter Source Map}

The table maps each topic in this chapter to its primary repository authority.
Detailed inventories remain in the subsystem chapters and appendices.

\small
\begin{longtable}{L{0.23\textwidth} L{0.69\textwidth}}
\caption{Authoritative sources for the project overview and boundaries}
\label{tab:orientation-source-map}\\
\toprule
\tableheader Topic & Authoritative repository path(s) \\
\midrule
\endfirsthead
\toprule
\tableheader Topic & Authoritative repository path(s) \\
\midrule
\endhead
Simulation profiles and composition &
  \file{mfja_3rd_floor_bringup/launch/room_315_only.launch.py};
  \file{mfja_3rd_floor_bringup/launch/full_floor.launch.py};
  \file{mfja_3rd_floor_bringup/launch/single_industrial_robot.launch.py};
  \file{mfja_3rd_floor_bringup/launch/room_315_floor_common.py};
  \file{mfja_robot_control_config/launch/multi_robot_sim.launch.py} \\ \rowseparator
Rail and robot simulation boundaries &
  \file{mfja_robot_control_config/scripts/room_315_kinematic_shuttle.py};
  \file{mfja_robot_control_config/scripts/room_315_kinematic_shuttle_node.py};
  \file{mfja_3rd_floor_description/src/SmoothJointTrajectoryController.cc};
  \file{mfja_3rd_floor_description/src/SymmetricGripperController.cc} \\ \rowseparator
Visual, language, and execution boundaries &
  \file{mfja_robot_control_config/scripts/room_315_visual_state_inference_node.py};
  \file{mfja_robot_control_config/scripts/room_315_visual_runtime_v4.py};
  \file{mfja_robot_control_config/scripts/room_315_task_goal_builder.py};
  \file{mfja_robot_control_config/scripts/room_315_task_execution.py};
  \file{mfja_robot_control_config/scripts/room_315_rail_safety_supervisor.py} \\ \rowseparator
Supported and optional dependency surfaces &
  \file{mfja_rail_interfaces/package.xml};
  \file{mfja_3rd_floor_description/package.xml};
  \file{mfja_robot_control_config/package.xml};
  \file{mfja_3rd_floor_bringup/package.xml};
  \file{mfja_robot_control_config/config/room_315_visual_state/visual_state_runtime.yaml};
  \file{mfja_robot_control_config/config/room_315_task_execution/task_execution_runtime.yaml};
  \file{mfja_robot_control_config/config/room_315_task_goal/task_goal_understanding.yaml};
  \file{mfja_robot_control_config/scripts/room_315_visual_state_to_lerobot.py} \\ \rowseparator
Onboarding procedure &
  \file{docs/INSTALLATION.md};
  \file{docs/QUICK_START_AND_FEATURE_GUIDE.md};
  \file{docs/TROUBLESHOOTING.md};
  \file{mfja_3rd_floor_bringup/launch/room_315_only.launch.py} \\ \rowseparator
Repository structure baseline &
  \file{mfja_rail_interfaces/CMakeLists.txt};
  \file{mfja_3rd_floor_description/CMakeLists.txt};
  \file{mfja_robot_control_config/CMakeLists.txt};
  \file{mfja_3rd_floor_bringup/CMakeLists.txt} \\
\bottomrule
\end{longtable}
\normalsize

\section{System Purpose and Capabilities}

The maintained project provides a focused Room 315 simulation, the complete
MFJA third-floor simulation, and an isolated industrial-robot simulation for
fast arm and gripper work. Chapter~\ref{ch:launch} identifies the stable launch
entry point and defaults for each profile.

Within Room 315, the software can spawn configured robots, command robot joints
and mobile bases through namespaced bridges, move multiple shuttles on
calibrated directed rail geometry, command switches and stoppers, publish
binary occupancy sensors, attach simulation-only payload visuals, and expose
the maintained control surface through typed ROS 2 interfaces.

An optional neuro-symbolic workflow adds paired-camera visual-state estimation,
constrained English task understanding, PDDL/PlanSys2 planning, one-step
closed-loop execution, deterministic rail-safety supervision, effect
verification, dataset generation, model training, artifact promotion, and
fault/acceptance campaigns. The language model and visual-state model are
separate and neither selects nor publishes actions. Enabling the
perception-and-safety branch does not implicitly start learned inference or task
execution.

\section{Supported and Optional Surfaces}

\begin{longtable}{L{0.24\textwidth} L{0.33\textwidth} L{0.35\textwidth}}
\caption{Supported runtime surfaces and prerequisites beyond a clean clone}
\label{tab:supported-runtime-surfaces}\\
\toprule
\tableheader Surface & Clean clone/base installation & Extra prerequisite \\
\midrule
\endfirsthead
\toprule
\tableheader Surface & Clean clone/base installation & Extra prerequisite \\
\midrule
\endhead
Worlds and robots & Available & Working OpenGL display only for GUI \\ \rowseparator
Kinematic rails and typed devices & Available & None beyond ROS/Gazebo dependencies \\ \rowseparator
Camera and primitive supervisor topics & Available, disabled by default & Enable the perception-and-safety launch branch \\ \rowseparator
Dataset generation tools & Code/config available & Disk space; some workflows need long Gazebo runs \\ \rowseparator
Learned visual inference & Code available, model absent & Isolated Torch environment and complete promoted V4 bundle \\ \rowseparator
English semantic model & Deterministic parser remains available & \code{llama-cpp-python} and separately downloaded GGUF checkpoint \\ \rowseparator
PlanSys2 execution & Code/config available, actuation disabled & PlanSys2 packages plus exact V4 promotion and task authorization artifacts \\ \rowseparator
LeRobot export & Export adapter available & Separately installed LeRobot stack \\
\bottomrule
\end{longtable}

\section{Boundaries and Non-goals}

\begin{longtable}{L{0.22\textwidth} L{0.35\textwidth} L{0.35\textwidth}}
\caption{Maintained project boundaries and explicitly excluded assurances}
\label{tab:orientation-boundaries}\\
\toprule
\tableheader Boundary & Inside this project & Explicitly outside \\
\midrule
\endfirsthead
\toprule
\tableheader Boundary & Inside this project & Explicitly outside \\
\midrule
\endhead
Simulation & Gazebo worlds, models, kinematic joints, pose services, cameras, and ROS bridges & Validated equivalence to every physical sensor, drive, PLC, and collision envelope \\ \rowseparator
Rail propulsion & Arc-length motion over calibrated directed CSV geometry & Wheel/rail contact, motor dynamics, traction, and certified block control \\ \rowseparator
Robot motion & Deterministic visual joint trajectories and symmetric gripper articulation & Force-controlled grasping, payload attachment, and validated collision avoidance \\ \rowseparator
Perception & Visual shuttle segment/position, bounding box, and loaded-state facts after strict acceptance & Direct learned planning, direct learned actuation, or controller pose used to improve visual localization \\ \rowseparator
Language & Constrained \code{TaskGoal} drafts, deterministic validation, clarification, and confirmation & Arbitrary English converted directly into robot programs \\ \rowseparator
Execution & Gazebo-only, checksum-authorized, one-step planning through a deterministic supervisor & Physical-machine authorization or bypass of supervisor/effect checks \\
\bottomrule
\end{longtable}

The complete deployment boundary and known-gap register are maintained in
Chapter~\ref{ch:handover}. Subsystem chapters state only the additional local
limitations needed to interpret their behavior safely.

\section{Source-of-Truth Order}

When two statements disagree, use this order:

\begin{enumerate}
  \item \file{mfja_rail_interfaces/msg/*.msg} and
        \file{mfja_rail_interfaces/srv/*.srv} for wire fields;
  \item current launch code and \code{ros2 launch ... --show-args} for launch
        inputs and defaults;
  \item versioned YAML, JSON, PDDL, CSV, SDF, and URDF for declarative contracts;
  \item node or plugin implementation for runtime behavior;
  \item maintained operator and architecture documentation for supported
        procedures; and
  \item historical reports, audits, presentations, and evidence for provenance
        only.
\end{enumerate}

Never copy a value from a generated \file{install/} tree back into source. A
symlink install can make source and install appear identical, but interfaces,
CMake metadata, compiled plugins, and many installed files still require a
rebuild.

\section{First Twenty Minutes for a New Maintainer}

\begin{procedurebox}[Onboarding checklist]
\begin{enumerate}
  \item Confirm the checked-out branch and commit with \code{git status} and
        \code{git log -1 --oneline}.
  \item Read the non-negotiable deployment boundary in
        Chapter~\ref{ch:handover} and the package ownership map in
        Chapter~\ref{ch:packages}.
  \item Create or identify the colcon workspace; keep this repository as one
        directory under \file{src/}.
  \item Follow the dependency and build procedure in
        Chapter~\ref{ch:installation}.
  \item Launch the headless Room 315 smoke profile from
        Chapter~\ref{ch:operations}.
  \item Verify \rosname{/clock}, the three world services, and one rail state
        and sensor topic.
  \item Run the smallest test that covers the intended first change.
  \item Before editing, identify the source owner using
        Chapter~\ref{ch:configuration} and the Repository Source Index in
        Appendix~\ref{app:catalog}.
\end{enumerate}
\end{procedurebox}
